\section{Represent arbitrary bytes with a restricted alphabet}
Base64 represents binary data using a small set of printable characters. Three input
bytes contain twenty-four bits; four output characters each represent six of those
bits. The transformation preserves the bytes, including zero bytes that would terminate
a C-style string. Its output usually occupies more space than its input. A recipient
can recover the original bytes from the alphabet and the encoded text without a secret.

Begin with three bytes you can trace on paper. Write a scalar encoder that separates
four six-bit values and uses them to select characters. Extend it to the final one or
two input bytes. Then reverse that calculation, checking the text before writing any
decoded output. Once the scalar interfaces agree with independent results, place four
complete groups in MSA byte lanes. Keep short final groups on a scalar path, where their
padding and exact stores remain explicit.

This project uses the standard alphabet with \texttt{+} and \texttt{/}, required final
padding, and no line wrapping. Its decoder accepts the canonical form: only alphabet
characters and correctly placed padding, with all unused bits zero. These choices fix
what counts as a valid input. Other protocols can choose different alphabets or handling
of whitespace; a decoder needs a named policy before its error cases can be tested.

The first complete observation is decimal bytes 77, 97 and 110 becoming the four
characters \texttt{TWFu}. Keep the numerical sextets visible before looking up the
characters. If those values are wrong, changing the alphabet table cannot repair the
bit split. If they agree but the characters do not, inspect the lookup or vector
translation rather than shifting the original bytes again.
